\section{Introducción}

HTCondor es un sistema de gestión de carga de trabajo especializado en computación distribuida de alto rendimiento (High Throughput Computing - HTC). Desarrollado originalmente en la Universidad de Wisconsin-Madison, HTCondor permite a organizaciones e instituciones académicas aprovechar recursos computacionales distribuidos para ejecutar grandes volúmenes de trabajos computacionales de manera eficiente y confiable.
A diferencia de otros sistemas de computación paralela que se enfocan en la velocidad de ejecución de trabajos individuales, HTCondor está diseñado para maximizar el rendimiento total del sistema, priorizando la capacidad de procesar miles o millones de trabajos a lo largo del tiempo. Esta característica lo convierte en una herramienta fundamental para investigadores y organizaciones que requieren procesar grandes cantidades de datos o ejecutar simulaciones extensivas.
El concepto de ``universo'' en HTCondor se refiere a diferentes entornos de ejecución que definen cómo y dónde se ejecutan los trabajos. Cada universo proporciona un conjunto específico de características y capacidades, desde la ejecución simple de programas hasta entornos virtualizados complejos. Los principales universos incluyen el universo estándar, vanilla, grid, docker, y virtual machine, cada uno optimizado para diferentes tipos de cargas de trabajo y requisitos de ejecución.
La importancia de HTCondor en el ecosistema de computación distribuida radica en su capacidad para democratizar el acceso a recursos computacionales de alto rendimiento. Mediante la implementación de políticas flexibles de asignación de recursos y un sistema robusto de tolerancia a fallos, HTCondor permite que organizaciones con recursos limitados puedan competir en investigación computacionalmente intensiva.
Este estudio de mapeo sistemático tiene como objetivo analizar el impacto de los diferentes universos de HTCondor en los dominios de computación distribuida y paralela, evaluando su efectividad, casos de uso, y contribuciones al avance del campo de la computación de alto rendimiento.